\begin{frame}{ReLU: 2012년 이후 은닉층의 사실상 표준}
  \begin{itemize}
    \item $f(z) = \max(0, z)$ — 미분이 $z > 0$에서 \textbf{정확히 1}이라
      곱해도 줄어지지 않음
    \item \textbf{계산이 가장 단순}: 지수함수 없음, 비교 연산 하나
    \item 음수 구간 출력이 정확히 0 $\Rightarrow$ \kb{희소(sparsity)} —
      활성된 뉴런만 신호 전달, 계산 빨라짐
    \item 2012년 크리즈베스키 일행: ``rectifier는 학습 시간을
      \textbf{크게 줄인다}'' — 같은 해 \kb{AlexNet}이 ReLU를 써
      ImageNet을 휩쓸며 은닉층 \textbf{사실상 표준}
  \end{itemize}
  \vskip0.4em
  \begin{block}{자주 하는 실수: 로지스틱회귀의 습관으로 은닉층에 시그모이드를}
    확률 해석이 필요한 것은 \textbf{출력층뿐}. 은닉층의 기본값은
    \kb{ReLU 계열}, 출력층은 \textbf{문제 형태에 맞게}:
    회귀 $\to$ 항등, 이진분류 $\to$ 시그모이드, 다중분류 $\to$ softmax.
  \end{block}
\end{frame}
